\begin{center}
\textbf{Resumo}
\end{center}

\noindent Este artigo tem como objetivo identificar como a literatura científica integra sustentabilidade, manutenção e gestão de edificações e métodos de apoio multicritério à decisão, e em que medida essa integração oferece bases para priorizar intervenções em edificações públicas universitárias. Para isso, foi conduzida uma revisão integrativa sistematizada, com apoio bibliométrico em três bases (Scopus, Web of Science e Crossref, 2010--2026) e síntese temática apoiada em título, resumo e palavras-chave, sem leitura de texto completo, resultando em um núcleo final de 104 registros após triagem, auditoria qualitativa e reavaliação criteriosa. Os resultados indicam predominância de critérios técnico-operacionais e institucionais de sustentabilidade sobre critérios ambientais, econômicos e sociais explicitamente nomeados, além de maior frequência de instrumentos genéricos de apoio à decisão (\textit{framework}, \textit{decision support}, BIM) do que de métodos multicritério formalmente nomeados (AHP, TOPSIS, MCDM), ainda minoritários no núcleo final. A aplicabilidade a instituições públicas de ensino superior permanece pouco explorada no corpus internacional sistematizado, com lacuna específica identificada em 11,5\% dos registros do núcleo final, o que sustenta a proposição de uma matriz analítica orientada a investigações futuras, não uma solução já consolidada pela literatura.

\medskip
\noindent\textbf{Palavras-chave:} manutenção predial; sustentabilidade; métodos multicritério de decisão; gestão de edificações públicas; ambiente construído.

\bigskip

\begin{center}
\textbf{Abstract}
\end{center}

\noindent This article aims to identify how the scientific literature integrates sustainability, building maintenance/management, and multi-criteria decision support methods, and to what extent this integration provides a basis for prioritizing interventions in public university buildings. A systematized integrative review was conducted, with bibliometric support across three databases (Scopus, Web of Science, and Crossref, 2010--2026) and thematic synthesis based on title, abstract, and keywords, without full-text reading, resulting in a final core of 104 records after screening, qualitative audit, and careful reassessment. Results indicate a predominance of technical-operational and institutional sustainability criteria over explicitly named environmental, economic, and social criteria, along with a higher frequency of generic decision-support instruments (\textit{framework}, \textit{decision support}, BIM) than formally named multi-criteria methods (AHP, TOPSIS, MCDM), still a minority in the final core. Applicability to public higher education institutions remains underexplored in the systematized international corpus, with a specific gap identified in 11.5\% of the final core records, supporting the proposal of an analytical matrix aimed at guiding future research rather than a solution already consolidated by the literature.

\medskip
\noindent\textbf{Keywords:} building maintenance; sustainability; multi-criteria decision-making; public building management; built environment.

\bigskip
